\SourceLeaf{182}\clearpage
\begin{center}
{\Large BRIEFWECHSEL.}
\end{center}

\bigskip
\begin{center}
[DIE TRANSCENDENTE TRIGONOMETRIE.]
\end{center}

\begin{center}\rule{2cm}{0.4pt}\end{center}

\begin{center}
[\textsc{Wachter} \emph{an} \textsc{Gauss}. \emph{Danzig, 12. December 1816.}]
\end{center}

\{.\ .\ .\ .\ .\ .\ .\ .\ Also die anti-Euklideische oder Ihre Geometrie wäre wahr.
Die \textls{Constante} in ihr aber bleibt unbestimmt: warum? liesse sich vielleicht durch
Folgendes begreiflich machen. .\ .\ .\ .\ .\ .\ .

.\ .\ .\ .\ .\ .\ .\ .\ Das Resultat des Bisherigen wäre also so auszusprechen:

Die
Euklideische Geometrie ist falsch; aber dennoch muss die wahre Geometrie mit
demselben elften Euklideischen Axiom oder mit dem Postulat von Linien und Flächen
anfangen, welche die in jenem Axiom behauptete Eigenschaft haben. Statt der
geraden Linie und Ebene sind nur die grössten Kreise jener mit unendlichem Radius
beschriebenen Kugel nebst ihrer Oberfläche zu setzen. Es entsteht zwar die eine
Unbequemlichkeit daraus, dass die Theile dieser Fläche bloss symmetrisch, nicht,
wie bei der Ebene, congruent sind; oder dass der Radius nach der einen Seite hin
unendlich, nach der anderen imaginär ist; allein wie jene Unbequemlichkeit durch
viele andere Vortheile, welche die Construction auf einer Kugelfläche darbietet,
wieder aufgewogen werde, ist klar: \textls{so dass vielleicht auch dann noch}, wenn die
Euklideische Geometrie wahr wäre, zwar nicht mehr die Nothwendigkeit obwaltete,
die Ebene als eine unendliche Kugelfläche zu betrachten, aber doch noch die
Fruchtbarkeit dieser Ansicht dieselbe empfehlen könnte.

Allein, als ich alles diess durchdacht, als ich mich über das Resultat schon
völlig beruhigt hatte, theils weil ich glaubte, der Grund (la métaphysique) jener
der Geometrie nothwendig anhaftenden Unbestimmtheit, --- auch
